\apendicemanual{Questionário de requisitos do Beto}{ap:levantamento-beto}

\begingroup\small
\noindent\textbf{Levantamento de Requisitos para Aprimoramento do Simulador Beto}

\textbf{Termo de consentimento}\par
Li as informações apresentadas acima e concordo em participar voluntariamente deste levantamento de requisitos, autorizando o uso anonimizado das minhas respostas para fins acadêmicos.
\begin{itemize}[label=$\circ$]
\item Concordo em participar do levantamento de requisitos.
\item Não concordo em participar do levantamento de requisitos.
\end{itemize}

\textbf{Coleta de e-mail}\par
E-mail: \rule{.8\linewidth}{.3pt}

\textbf{Perfil: familiaridade}\par
Qual a sua familiaridade com o conteúdo de projeto de memória cache (normalmente visto nas disciplinas de AOC1 e AOC2)?
\begin{itemize}[label=$\circ$]
\item Nenhuma
\item Básica
\item Intermediária
\item Avançada
\end{itemize}

\textbf{Perfil: disciplinas}\par
Você está cursando Arquitetura e Organização de Computadores II, seja laboratório ou teórica, ou é ex-aluno(a) da(s) disciplina(s)?
\begin{itemize}[label=$\circ$]
\item Cursando somente AOCII
\item Cursando somente LAOCII
\item Cursando AOCII e LAOCII
\item Ex-aluno AOCII e/ou LAOCII
\end{itemize}

\Needspace{17\baselineskip}
\textbf{Pergunta 1 --- Configuração da simulação}\par
Durante a configuração de uma simulação no Beto, quais dificuldades você encontrou?

\emph{Múltipla seleção.}
\begin{itemize}[label=\caixaOpcao]
\item Entender o significado dos campos de configuração
\item Saber quais valores eram válidos para os campos
\item Escolher o tipo de mapeamento
\item Escolher a política de substituição
\item Entender a relação entre memória principal, bloco e linhas de cache
\item Identificar quais campos precisavam ser preenchidos antes de iniciar a simulação
\item Campos desalinhados ou visualmente confusos
\item Falta de mensagens de erro claras quando algum valor estava incorreto
\item Não encontrei dificuldades nessa etapa
\item Outro: \rule{.65\linewidth}{.3pt}
\end{itemize}

\Needspace{17\baselineskip}
\textbf{Pergunta 2 --- Entendimento da execução}\par
Durante a simulação no Beto, quais informações você sentiu falta para entender melhor como cada endereço acessado era tratado pela memória cache?

\emph{Múltipla seleção.}
\begin{itemize}[label=\caixaOpcao]
\item Campo de tag do endereço
\item Índice ou conjunto acessado
\item Deslocamento dentro do bloco
\item Indicação mais clara de acerto (hit) ou falta (miss)
\item Linha ou via escolhida na cache
\item Motivo da substituição realizada
\item Valor ou estado LRU de cada linha ou via
\item Bit de validade
\item Dirty bit/bit de sujeira
\item Histórico passo a passo dos acessos realizados
\item Estatísticas finais, como quantidade de hits, misses e taxa de acerto
\item Não senti falta de informações nessa etapa
\item Outro: \rule{.65\linewidth}{.3pt}
\end{itemize}

\Needspace{17\baselineskip}
\textbf{Pergunta 3 --- Problemas de usabilidade/UX}\par
Quais problemas de usabilidade ou experiência de uso você percebeu ao usar o Beto?

\emph{Múltipla seleção.}
\begin{itemize}[label=\caixaOpcao]
\item Falta de feedback após executar uma ação
\item Mensagens pouco claras
\item Termos técnicos sem explicação suficiente
\item Dificuldade para encontrar botões ou controles
\item Dificuldade para corrigir um erro durante a configuração ou execução
\item Inconsistência visual entre telas/campos
\item Excesso de informação na tela
\item Desalinhamento de campos ou rótulos
\item Falta de ajuda, instruções ou documentação dentro da interface
\item Não percebi problemas de usabilidade
\item Outro: \rule{.65\linewidth}{.3pt}
\end{itemize}

\Needspace{17\baselineskip}
\textbf{Pergunta 4 --- Priorização de melhorias}\par
Considerando as dificuldades e problemas percebidos durante o uso do Beto, quais melhorias você considera mais importantes para uma nova versão da ferramenta? (Marque todas as opções que achar você achar ser essencial para compor uma ferramenta com o propósito do Beto)

\emph{Múltipla seleção.}
\begin{itemize}[label=\caixaOpcao]
\item Melhorar a validação de entradas antes de iniciar a simulação
\item Apresentar mensagens de erro mais claras e didáticas
\item Reorganizar visualmente os campos de configuração
\item Explicar melhor os termos técnicos usados na interface
\item Mostrar a decomposição do endereço em tag, índice/conjunto e deslocamento
\item Exibir uma tabela de histórico com cada acesso executado
\item Mostrar de forma mais clara quando ocorreu hit ou miss
\item Mostrar a linha ou via escolhida em cada acesso
\item Exibir o valor ou estado LRU de cada linha ou via
\item Exibir o bit de validade
\item Incluir suporte a operações de leitura e escrita (load/store)
\item Incluir suporte a dirty bit e políticas de escrita
\item Apresentar estatísticas finais da simulação, como hits, misses e taxa de acerto
\item Disponibilizar exemplos prontos de simulação
\item Melhorar a ajuda ou documentação dentro da própria interface
\item Outro: \rule{.65\linewidth}{.3pt}
\end{itemize}

\Needspace{7\baselineskip}
\textbf{Pergunta 5 --- Campo aberto}\par
Além das opções marcadas anteriormente, há alguma melhoria que você considera importante para uma nova versão do Beto? Se sim, descreva brevemente.
\par\vspace{.5em}\noindent\rule{\linewidth}{.3pt}\par\vspace{1em}\noindent\rule{\linewidth}{.3pt}

\Needspace{6\baselineskip}
\textbf{Pergunta 6 --- Avaliação geral}\par
De 1 a 4, quão útil você considera o Beto para apoiar o entendimento de memória cache?
\par\noindent Seleção de um valor: $\circ$ 1\qquad $\circ$ 2\qquad $\circ$ 3\qquad $\circ$ 4.
\endgroup
